\section{Conclusion}
\label{sec:conclusion}

Agent Manifest gives an AI agent a verifiable identity. By binding the ten artifacts
that define an agent (system prompt, policy, tools, model, retrieval corpus,
memory, decision trace, delegation, supply chain, and human approvals) into one
signed, hardware-attestable, tamper-evident record, it lets a verifier detect differences between supplied artifacts and a signed
baseline. Establishing that those artifacts were used in a deployed agent requires
additional runtime evidence and trusted attestation binding. The design favors
tamper-evidence over tamper-resistance, reuses established primitives (SPIFFE, JCS,
Ed25519/ML-DSA, RFC~9162 Merkle trees, RATS/EAT), and slots into the runtime stack
alongside cMCP and TRACE for attestation and AGT for policy.

\paragraph{Limitations and future work.}
The format is settled for the static artifacts; several pieces are staged. Memory and
retrieval bindings capture a baseline and a drift policy but continuous in-flight
verification of mutating state is future work. The post-quantum profile is specified
and implemented but was not exercised: its tests skip without the ML-DSA
dependency. Federated, cross-organization
verification (where one trust domain appraises another's manifests) is targeted
for the next revision, alongside richer multi-model manifests. The specification is a
developer-preview draft; the roadmap is community hardening toward an open standard
with multiple independent implementations and conformance certification.
